\part{Algebraic Topology}
\chapter{Homotopy}
\begin{importantdefinition}{}{}
	Let $X,Y$ be topological spaces. Let $f,g : X \to Y$ be continuous. Then $f$ is called \tib{homotopic} to $g$ via a map $h$ iff. $h$ is a "continous interpolation between $f$ and $g$, i.e:
	\begin{enumerate}
		\item $h$ is a continous map $X \times [0,1] \to Y$,
		\item For all $x \in X$, we have $h(x,0) = f(x)$,
		\item For all $x \in X$, we have $h(x,1) = g(x)$.
	\end{enumerate}
	We call $h$ a \tib{free homotopy}, or simply a \tib{homotopy}, between $f$ and $g$. We write $f \sim_h g$ if $f$ is a homotopy between $f$ and $g$. We write $f \sim g$ if there exists a homotopy between $f$ and $g$.
\end{importantdefinition}
\begin{exercise}
	Show that the property of being homotopic is an equivalence relation.
\end{exercise}
\begin{definition}
	We say that $f$ and $g$ are \tib{homotopic relative to a subset $A \subseteq X$} and write $f \sim_{\rel A}$ or $f \sim_{h, \rel A}$ iff:
	\begin{enumerate}
		\item $f \sim_h g$,
		\item For all $t \in [0,1]$, we have $h(t, \cdot)|_A = f|_A = g|_A$
	\end{enumerate}
\end{definition}
I personally find the designation "homotopic relative to $A$" quite unintuitive, given that the property of being homotopic relative to a given subset is stronger than the property of being homotopic.
\begin{exercise}
	Show that $f \sim_h g$ is equivalent to $f \sim_{h, \rel \emptyset} g$.
\end{exercise}
\begin{exercise}
	Let $f_0, f_1 : X \to Y$ and $g_0, g_1 : Y \to Z$ be continuous. Let $h_0$ be a homotopy between $f_0$ and $f_1$ and let $h_1$ be a homotopy between $f_1$ and $g_1$. Then $H : (x,t) \mapsto h_1(h_0(x,t),t)$ is a homotopy between $g_0 \circ f_0$ and $g_1 \circ f_1$.
\end{exercise}

\begin{importantdefinition}{}{}
	Let $f : X \to Y$ be continuous. Then $f$ is called a \tib{homotopy equivalence} between $X$ and $Y$ iff. there exists a map $g : Y \to X$ such that $g \circ f$ is homotopic to $\id_X$ and $f \circ g$ is homotopic to $\id_Y$. We write $X \simeq_{f,g} Y$, $X \simeq_f Y$, or simply $X \simeq Y$.
\end{importantdefinition}
\begin{exercise}
	Show that the property of being homotopy equivalent is an equivalence relation.
\end{exercise}
\begin{exercise}
	Show that homeomorphism implies homotopy equivalence.
\end{exercise}
\begin{definition}
	A topological space $X$ is called \tib{contractible} iff. $X$ is homotopy equivalent to a one-point space.
\end{definition}
\begin{exercise}
	\begin{enumerate}
		\item Show that
		\begin{align*}
			h : (x,t) \mapsto x - tx
		\end{align*}
		is a homotopy between $\id_{\bR^n}$ and $f : \bR^n \mapsto \lr{0}$.
		\item Conclude that $\bR^n$ is contractible.
		\item Conclude that homotopy equivalence does not imply homeomorphism.
	\end{enumerate}
\end{exercise}
\begin{exercise}
	A topological space $X$ is contractible if and only if there exists a point $x_0 \in X$ such that $\id_X$ is homotopic to the constant map $X \to {x_0}$.
\end{exercise}
\begin{definition}
	Let $X$ be a subset of a topological vector space. Then we call $X$ \tib{star-shaped} iff. there exists a point $x_0 \in X$ such that $X$ contains the connecting line between $x_0$ and any other point $x \in X$, i.e.
	\begin{align*}
		\forall x \in X:\\
		\forall t \in [0,1]:\\	
		(1-t)x_0 + tx \in X
	\end{align*}
\end{definition}
\begin{exercise}
	Show that any star-shaped set is contractible. Conclude that any convex set is contractible.
\end{exercise}
\begin{lemma}
	Let $X$ be a topological space. Let $A \subseteq X$. Then $A$ and $X$ are homotopy equivalent if there exists a homotopy between $\id_X$ and $\id_X|_A$, i.e. a continuous map $h : X \times [0,1] \to X$ such that for all $x \in X$:
	\begin{enumerate}
		\item $h(x,0) = x$,
		\item $h(x,1) \in A$,
		\item $x \in A \implies h(x,1) = x$
	\end{enumerate}
\end{lemma}

\clearpage

\section{The Fundamental Group}
\begin{definition}
	Let $X$ be a topological space. Then a continuous map $\gamma [0,1] \to X$ with $\gamma(0) = \gamma(1)$ is known as a \tib{loop}. We denote the set of all loops in $X$ with $\gamma(0) = \gamma(1) = x_0$ as $\Omega(X, x_0)$.
\end{definition}
\begin{lemma}
	Let $\gamma : [0,1] \to X$ be a loop. Let $\phi : [0,1] \to [0,1]$ be a continuous map such that $\gamma(0) = 0$ and $\gamma(1) = 1$. Then $\gamma \circ \phi$ and $\gamma$ are homotopic relative to $\lr{0,1}$.
\end{lemma}
\begin{proof}
	The map
	\begin{align*}
		h : (x,t) \mapsto \gamma((1-t)x + t \phi(x))
	\end{align*}
	is a homotopy between $\gamma$ and $\gamma \circ \phi$.
\end{proof}
\begin{definition}
	Let $f, g \in \Omega(X, x_0)$. Then we define the \tib{concatenation of $f$ and $g$} as the map
	\begin{align*}
		(f * g)(t) =
		\begin{cases}
			f(2t) & 0 \leq t \leq 1/2,\\
			g(2t - 1) & 1/2 \leq t \leq 1.
		\end{cases}
	\end{align*}
\end{definition}
\begin{definition}
	Let $\gamma \in \Omega(X, x_0)$. We define the \tib{inverse loop} of $\gamma$ to be the loop that "moves through gamma backwards", i.e.
	\begin{align*}
		\ol \gamma(t) := \gamma(1-t)
	\end{align*}
\end{definition} 
\begin{exercise}
	Verify that $\ol{\ol{\gamma}} = \gamma$.
\end{exercise}
\begin{exercise}
	Let $\gamma \in \Omega(X, x_0)$. Then $\gamma * \ol \gamma$ is homotopic to $e_{x_0} : X \to \lr{x_0}$ relative to $\lr{0,1}$.
\end{exercise}
\begin{importanttheorem}{Existence of the fundamental group}{}
	Let $X$ be a topological space. Let $x_0 \in X$. Then the quotient space
	\begin{align*}
		\pi_1(X, x_0) = \Omega(X, x_0) / \sim_{\rel({0,1})}
	\end{align*}
	of loops modulo homotopy
	becomes a group, known as the \tib{fundamental group of $(X, x_0)$}, when equipped path concatenation
	\begin{align*}
		[f] * [g] = [f * g],
	\end{align*}
	as the group operation.
	The neutral element of this group is $e_{x_0} := (x \mapsto x_0)$.
\end{importanttheorem}
\begin{theorem}
	If we additionally define
	\begin{align*}
		\pi_1(f) : \pi_1(X,x) &\mapsto \pi_1(Y,f(x))\\
		[\gamma] &\mapsto [f \circ \gamma]
	\end{align*}
	for all continuous $f : X \to Y$, 
	the map $\pi_1$ becomes a covariant functor from the category of topological spaces into the category of groups.
\end{theorem}
\begin{definition}
	Whenever we need to make the base point clear, we write $\pi_1(f,x) := \pi_1(f) : \pi_1(X,x) \mapsto \pi_1(Y,f(x))$
\end{definition}
\begin{importanttheorem}{Fundamental groups at different base points}{}
	Let $X$ be a topological space. Let $x,y \in X$. Then every continuous path $\delta : [0,1] \to X$ induces a group isomorphism
	\begin{align*}
		J_\delta : \pi_1(X, \delta(0)) &\cong \pi_1(X,\delta(1))\\
		[\gamma] &\mapsto \lr[\ol \delta * \gamma * \delta]
	\end{align*}
\end{importanttheorem}
\begin{proof}
	\begin{enumerate}
		\item Since loop concatenation preserves homotopy, $J_\delta$ sends different representatives of the same equivalence class to the same equivalence class.
		\item $J_\delta$ goes from $\pi_1(X,x)$ to $\pi_1(X,y)$, since for any $\gamma : [0,1] \to X$ with $\gamma(0) = \gamma(1) = x$, $\ol \delta$ traces out a path from $\delta(1)$ to $\delta(0)$, then $\gamma$ traces out a loop from $\delta(0)$ to $\delta(0)$, and then $\delta$ traces out a path from $\delta(0)$ back to $\delta(1)$, i.e. $\ol \delta * \gamma * \delta$ is a loop from $\delta(1)$ to $\delta(1)$.
		\item $J_\delta$ is a group homomorphism, since
		\begin{align*}
			J_\delta([\gamma_1 * \gamma_2])
			&=
			\lr[\ol \delta * \gamma_1 * \gamma_2 * \delta]\\
			&=
			\lr[\ol \delta * \gamma_1 * e_{x_0} * \gamma_2 * \delta]\\
			&=
			\lr[\ol \delta * \gamma_1 * \delta * \ol \delta * \gamma_2 * \delta]\\
			&=
			\lr[\ol \delta * \gamma_1 * \delta] * \lr[\ol \delta * \gamma_2 * \delta]\\
			&=
			J_\delta([\gamma_1]) * J_\delta([\gamma_2])
		\end{align*}
		(Recall that any map that respects the group operation also respects the neutral element and inverses, and is thus a homomorphism.)
		\item $J_\delta$ is bijective with $J_\delta ^{-1} = J_{\ol \delta}$:
		\begin{align*}
			J_\delta(J_{\ol \delta}(\gamma))
			&=
			\lr[\ol \delta * \lr(\delta * \gamma * \ol \delta) * \delta]
			=
			\lr[\gamma]\\
			J_{\ol \delta}(J_{\delta}(\gamma))
			&=
			\lr[\delta * \lr(\ol \delta * \gamma * \delta) * \ol \delta]
			=
			\lr[\gamma]
		\end{align*}
	\end{enumerate}
\end{proof}
\begin{corollary}
	If $X$ is path connected, $\pi_1(X,x) \cong \pi_1(X,y)$ for any $x,y \in X$.
\end{corollary}
\begin{definition}
	Let $\gamma_0 \in \Omega(X,x)$ an $\gamma_1 \in \Omega(X,y)$. Then we call $h : [0,1] \times [0,1] \to X$ a \tib{free homotopy of loops}, or simply \tib{loop homotopy}, between $\gamma_0$ and $\gamma_1$ iff. $h$ is a homotopy between $\gamma_0$ and $\gamma_1$ such that $h_t(x) := h(x,t)$ is a loop for any fixed $t \in [0,1]$.
\end{definition}
\begin{exercise}
	\begin{enumerate}
		\item Verify that $h$ is a loop homotopy iff. $h$ is a homotopy and $h_t(0) = h_t(1)$ for any $t \in [0,1]$.
		\item Show that any $\lr{0,1}$-homotopy is a loop homotopy.
	\end{enumerate}	
\end{exercise}
The additional constraint that and $h(\cdot, t)$ is a loop may seem superfluous, but there are in fact homotopies between loops that aren't loop homotopies:
\begin{exercise}
	\begin{enumerate}
		\item Let $\gamma_0 = \gamma_1 : [0,1] \to \lr{0}$ be constant loops. Let $h(x,t) := xt(1-t)$. Then $h$ is a homotopy between $\gamma_0$ and $\gamma_1$, but fails to be a loop homotopy.
		\item Let $\gamma_0(x) = e^{2\pi i x}$. Let $\gamma_1(x) = e^{4 \pi i x}$. Then $h(x,t) := e^{2 \pi i (1+t)x}$ is a homotopy between $\gamma_0$ and $\gamma_1$, but fails to be a loop homotopy.
		\item Assume that any loop $\gamma : [0,1] \to S^1$ with $1$ as its base point is homotopic a loop of the form $\gamma(x) := e^{2\pi n ix}$ for some $n \in \bZ$ (we will prove this later). Does 2. imply that the fundamental group $\pi_1(S^1,1)$ is trivial?
	\end{enumerate}
\end{exercise}
\begin{definition}
	A loop $\gamma \in \Omega(X,x)$ is called \tib{contractible in $X$} iff. there exists an $x \in X$ and a loop homotopy $h$ such that $\gamma \sim_h e_x$.
\end{definition}
\begin{lemma}
	$\gamma \in \Omega(X,x)$ is contractible in $X$ if $\gamma \sim_{\rel {0,1}} e_{\gamma(0)}$.
\end{lemma}
\begin{theorem}
	Let $\gamma \in \Omega(X, x)$ be a loop. Let $S^1 \cong I/\lr{0,1}$ be equipped with the quotient topology. Then $\gamma$ is contractible if and only if $\gamma \circ p : S^1 \to X$ has a continuous extension to the closed unit ball $\ol B^2$.
\end{theorem}
\begin{theorem}
	Let $h$ be a homotopy between $f,g : X \to Y$. Let $x_1 \in X$. Then
	\begin{align*}
		h(x_1,\cdot) : [0,1] &\to Y\\
		t &\mapsto h(x_1,t)
	\end{align*}
	is a path from $f(x_1)$ to $g(x_1)$
	which fulfills
	\begin{align*}
		\pi_1(g,x_1) = J_{h(x_1, \cdot)} \circ \pi_1(f, x_1)
	\end{align*}
\end{theorem}
Recall that $\pi_{1}(g, x_1)$ is the map $[\gamma] \mapsto [g \circ \gamma]$, and that $J_\delta$ is the map $[\gamma] \mapsto [\ol \delta * \gamma * \delta]$, so this theorem states that
\begin{align*}
	[g \circ \gamma] = \lr[\ol h(x_1,\cdot) * (f \circ \gamma) * h(x_1,\cdot)]
\end{align*}
for any $\gamma \in \pi_1(X,x_1)$, or, equivalently, that the diagram
% https://q.uiver.app/#q=WzAsMyxbMCwwLCJcXHBpXzEoWCx4XzEpIl0sWzIsMCwiXFxwaV8xKFksZih4XzEpKSJdLFsxLDIsIlxccGlfMShZLGcoeF8xKSJdLFsxLDIsIkpfe2goeF8xLFxcY2RvdCl9Il0sWzAsMiwiXFxwaV97MSx4XzF9KGcpIiwyXSxbMCwxLCJcXHBpX3sxLHhfMX0oZikiXV0=
\[\begin{tikzcd}
	& {\pi_1(X,x_1)} &\\
	\\
	{\pi_1(Y,f(x_1))} &&  {\pi_1(Y,g(x_1)}
	\arrow["{\pi_{1}(f,x_1)}", from=1-2, to=3-1]
	\arrow["{\pi_{1}(g,x_1)}"', from=1-2, to=3-3]
	\arrow["{J_{h(x_1,\cdot)}}",  from=3-1, to=3-3]
\end{tikzcd}\]
commutes.
\begin{proof}
	Define
	\begin{align*}
		H_t(x) := h(x_1, t + (1-t)x).
	\end{align*}
	We have $(f \circ \gamma) \in \Omega(X, f(x_1)$ and $(g \circ \gamma) \in \Omega(X, g(x_1))$, since
	\begin{align*}
		f(\gamma(0)) = f(x_1) = f(\gamma(1))\\
		g(\gamma(0)) = g(x_1) = g(\gamma(1))
	\end{align*}
	Then a homotopy between $g \circ \gamma$ and $\ol h(x_1,\cdot) * (f \circ \gamma) * h(x_1,\cdot)$ is given by:
	\begin{align*}
		k(x,t) := (\ol{H_t} * (h(\cdot, t) \circ \gamma) * H_t)(x),
	\end{align*}
	since we get:
	\begin{align*}
		k(x,0)
		&=
		(\ol{H_0} * (h(\cdot, 0) \circ \gamma) * H_0)(x)\\
		&=
		(\ol{H_0} * (f \circ \gamma) * H_0)(x)\\
		&=
		(\ol h(x_1, \cdot) * (f \circ \gamma) * h(x_1, \cdot))(x)
	\end{align*}
	and
	\begin{align*}
		k(x,1)
		&=
		(\ol{H_1} * (h(\cdot, 1) \circ \gamma) * H_1)(x)\\
		&=
		({h(x_1,(1-1) + (1-(1-1))(\cdot))} * (g \circ \gamma) * h(x_1,1 + (1-1)(\cdot)))(x)\\
		&=
		({h(x_1,\cdot)} * (g \circ \gamma) * h(x_1,1))(x)\\
		&=
		({e_{g(x_1)}} * (g \circ \gamma) * e_{g(x_1)})(x)\\
		&=
		(g \circ \gamma)(x),
	\end{align*}
	where the second to last step follows since the loop $(g \circ \gamma)$ ends on the base point $g(x_1)$. 
\end{proof}
\begin{lemma}
	Let $\gamma : [0,1] \to X$ be continuous. Let $f : X \to Y$. Then
	\begin{align*}
		\pi_{1}(f, \gamma(1)) \circ J_\gamma = J_{f \circ \gamma} \circ \pi_{1}(f, \gamma(0)),
	\end{align*}
	i.e. the diagram
	% https://q.uiver.app/#q=WzAsNCxbMCwwLCJcXHBpXzEoWCwgXFxnYW1tYSgxKSkiXSxbMCwyLCJcXHBpXzEoWCwgXFxnYW1tYSgwKSkiXSxbMiwwLCJcXHBpXzEoWSwgZihcXGdhbW1hKDEpKSkiXSxbMiwyLCJcXHBpXzEoWSwgZihcXGdhbW1hKDApKSkiXSxbMCwxLCJKX1xcZ2FtbWEiXSxbMiwzLCJKX3tmIFxcY2lyYyBcXGdhbW1hfSJdLFswLDIsIlxccGlfezEsIFxcZ2FtbWEoMSl9KGYpIiwxXSxbMSwzLCJcXHBpX3sxLCBcXGdhbW1hKDApfShmKSIsMV1d
	\[\begin{tikzcd}
		{\pi_1(X, \gamma(0))} && {\pi_1(Y, f(\gamma(0)))} \\
		\\
		{\pi_1(X, \gamma(1))} && {\pi_1(Y, f(\gamma(1)))}
		\arrow["{\pi_{1}(f, \gamma(0))}"{description}, from=1-1, to=1-3]
		\arrow["{J_\gamma}", from=1-1, to=3-1]
		\arrow["{J_{f \circ \gamma}}", from=1-3, to=3-3]
		\arrow["{\pi_{1}(f, \gamma(1))}"{description}, from=3-1, to=3-3]
	\end{tikzcd}\]
	commutes.
\end{lemma}
\begin{proof}
	Let $\sigma \in \Omega(X, \gamma(0))$. Then we have
	\begin{align*}
		\pi_{1}(f, \gamma(1)) \circ J_\gamma([\sigma])
		&=
		\pi_{1}(f, \gamma(1))(\lr[\ol \gamma * \sigma * \gamma])\\
		&=
		\lr[f \circ (\ol \gamma * \sigma * \gamma)]\\
		&=
		\lr[(f \circ \ol \gamma) * (f \circ \sigma) * (f \circ \gamma)]\\
		&=
		\lr[\ol{(f \circ \gamma)} * (f \circ \sigma) * (f \circ \gamma)]\\
		&=
		J_{f \circ \gamma}([f \circ \sigma])\\
		&=
		J_{f \circ \gamma} \circ \pi_{1}(f, \gamma(0))([\sigma])
	\end{align*}
\end{proof}
\begin{theorem}
	Let $f : X \to Y$ be a homotopy equivalence. Let $x_0 \in X$. Then the map $\pi_{1}(f, x_0)$ is a group isomorphism $\pi_1(X, x_0) \cong \pi_1(Y, f(x_0))$.
\end{theorem}
\chapter{Fibrations and Coverings}
\begin{importantdefinition}{Locally Trivial Fibration}{}
	Let $\pi : Y \to X$ be continuous. Then $\pi$ is called a \tib{locally trivial fibration} iff. for every $x \in X$, there exists an open neighborhood $U$ such that there exists a homeomorphism $h_u : \pi^{-1}[U] \cong U \times \pi^{-1}[x]$ fulfilling $p_1 \circ h_U = \pi|_{\pi^{-1}[U]}$.
\end{importantdefinition}
% https://q.uiver.app/#q=WzAsMyxbMCwwLCJZIFxcc3Vwc2V0ZXFcXHBpXnstMX1bVV0iXSxbMCwyLCJYIFxcc3Vwc2V0ZXEgVSJdLFsyLDAsIlUgXFx0aW1lcyBcXHBpXnstMX1beF0iXSxbMCwxLCJcXHBpIl0sWzAsMiwiaF9VIiwyXSxbMiwxLCJwXzEiXV0=
\[\begin{tikzcd}
	{Y \supseteq\pi^{-1}[U]} && {U \times \pi^{-1}[x]} \\
	\\
	{X \supseteq U}
	\arrow["{h_U}"', from=1-1, to=1-3]
	\arrow["\pi", from=1-1, to=3-1]
	\arrow["{p_1}", from=1-3, to=3-1]
\end{tikzcd}\]
\begin{definition}
	A locally trivial fibration is called a \tib{local covering} if every $\pi^{-1}[x]$ is discrete. A locally trivial fibration is called a \tib{trivial fibration} if we can always pick $U = X$, i.e. $h : \pi^{-1}[X] \cong X \times \pi^{-1}[x]$. A trivial fibration that is also a local covering is simply called a \tib{covering}.
\end{definition}
\begin{exercise}
	If we replace $\pi^{-1}[x]$ in the definition by an arbitrary topological space $F_U$, we get $F_U \cong \pi^{-1}[x]$. Thus, both definitions are equivalent.
\end{exercise}
\begin{lemma}
	Let $\pi : Y \to X$ be locally trivial fibrations. Let $X$ be connected. Then fibers of $\pi$ are homeomorphic.
\end{lemma}
\begin{theorem}
	Let $\pi : Y \to X$ be continuous. Then the following are equivalent:
	\begin{enumerate}
		\item $\pi$ is a local covering.
		\item For every $x \in X$, there exists an open neighborhood $U$ of $x$ and a non-empty $I$ such that for ecah $i \in I$, there exists an open set $V_i$ such that:
		\begin{enumerate}
			\item $\pi^{-1}[U] = \bigcup_{i \in I} V_i$.
			\item 
			For $i \neq j \in I$, we have $V_i \cap V_j = \emptyset$.
			\item
			For all $i \in I$, the restriction of $\pi$ to $V_i$ is a homeomorphism between $V_i$ and $U$.
		\end{enumerate}
	\end{enumerate}
\end{theorem}
\begin{importantdefinition}{Lifting}{}
	Let $\pi : Y \to X$ be a locally trivial fibration. Let $f : Z \to X$ be continuous. Then a continuous map $\hat f : Z \to Y$ is called a \tib{lift from $f$ to $Y$} if $\pi \circ \hat f = f$.
\end{importantdefinition}
Any path can be lifted:
\begin{theorem}
Let $\pi : Y \to X$ be a locally trivial fibration. Let $\gamma : [0,1] \to XX$ be continuous. Let $\pi(y_0) = \gamma(0)$. Then there exists a lift $\hat \gamma : [0,y] \to Y$ with $\hat \gamma(0) = y_0$ such that $\pi \circ \hat f = f$.
\end{theorem}
And lifts respect fundamental groups:
\begin{theorem}
	Let $\pi : Y \to X$ be a locally trivial fibrations. Let $f : Z \to X$. Let $\hat f : Z \to Y$ be a lift of $f$. Then we have:
	\begin{align*}
		\forall &z_0 \in Z:\\
		\exists &y_0 \in Y:\\
		&\hat f(z_0) = y_0,\\
		&\pi_1(f)[\pi_1(Z, z_0)] \subseteq \pi_1(\pi)[\pi_1(Y,y_0)]
	\end{align*}
\end{theorem}
\begin{proof}
	\begin{align*}
		\pi_1(f)[\pi_1(Z, z_0)]
		&=
		(\pi_1(\pi \circ \hat f))[\pi_1(Z, z_0)]\\
		&=
		(\pi_1(\pi) \circ \pi_1(\hat f))[\pi_1(Z, z_0)]\\
		&=
	\end{align*}
\end{proof}